%===============================================================================
% LLMWiki Technical Specification - Volume 01: Architecture
%===============================================================================
\chapter{Architecture}
\label{cha:architecture}

LLMWiki is not a Retrieval-Augmented Generation (RAG) framework. It is a \textbf{semantic operating layer} that continuously transforms heterogeneous digital artifacts into an evidence-aware knowledge space. This volume establishes the architectural foundations: the three-layer decomposition, the filesystem-as-WAL principle, the data flow invariants, and the module interface contracts that all subsequent volumes build upon.

\section{Design Philosophy}
\label{sec:architecture:philosophy}

\begin{designprinciple}[Graph as Source of Truth]
The property knowledge graph is the canonical knowledge representation. Vector indices are derived projections optimized for similarity search---not the primary store.
\end{designprinciple}

\begin{designprinciple}[Filesystem as Write-Ahead Log]
The local filesystem (directory structure .llmwiki/) serves as the sole durability mechanism. No auxiliary transaction log service (Kafka, Redis Streams, etc.) is required.
\end{designprinciple}

\begin{designprinciple}[Evidence Gate as First-Class Module]
Retrieval results must pass a mandatory cross-encoder verification stage before synthesis. This is not a reranker; it is an evidence filter with calibrated confidence thresholds.
\end{designprinciple}

\begin{designprinciple}[Local-First, LLM-Optional]
All critical-path operations (parsing, chunking, embedding, graph update, indexing, evidence verification) run locally with zero external API dependencies. LLMs are used only for planning, entity extraction (optional), and answer synthesis.
\end{designprinciple}

\begin{designprinciple}[Content Addressability]
Every semantic unit is identified by SHA-256 of its canonical form. Identical content implies identical identifier across all indices, enabling deduplication, caching, and verifiable lineage.
\end{designprinciple}

\begin{designprinciple}[Incremental Consistency]
After any ingestion, the system state (G, I, U) - graph, index, units - satisfies: a query executed against the live state produces the same result as replaying the ingestion log.
\end{designprinciple}

\begin{designprinciple}[Deterministic Replay]
Given identical .llmwiki/ directory and identical query q, two executions produce bitwise-identical results (modulo non-deterministic LLM sampling, which is confined to planning/synthesis).
\end{designprinciple}

\section{Three-Layer Architecture}
\label{sec:architecture:three-layer}

The system decomposes into three layers with strictly unidirectional data flow: Ingestion to Core to Query.

\begin{figure}[htbp]
\centering
\begin{adjustbox}{max width=\textwidth}
\begin{tikzpicture}[
  module/.style={draw, rounded corners, fill=white, text width=3.2cm, minimum height=0.8cm, font=\small, align=center, inner sep=3pt},
  layertitle/.style={font=\bfseries\small},
  arrow/.style={->, thick, >=stealth, shorten >=2pt, shorten <=2pt}
]
  % Layer 1: Ingestion
  \node[layertitle] at (0,3.3) {Layer 1: Ingestion};
  \node[module] (parse) at (0,2.5) {Semantic Parser (Vol. 02)};
  \node[module] (chunk) at (0,1.5) {Chunker (Vol. 02)};
  \node[module] (embed) at (0,0.5) {Embedder (Vol. 02)};

  % Layer 2: Core
  \node[layertitle] at (4.6,3.3) {Layer 2: Core};
  \node[module] (graph) at (4.6,2.5) {Knowledge Graph (Vol. 03)};
  \node[module] (index) at (4.6,1.5) {Semantic Index (Vol. 03)};
  \node[module] (evidence) at (4.6,0.5) {Evidence Gate (Vol. 06)};

  % Layer 3: Query
  \node[layertitle] at (9.2,3.3) {Layer 3: Query};
  \node[module] (planner) at (9.2,2.5) {Query Planner (Vol. 07)};
  \node[module] (runtime) at (9.2,1.5) {Agent Runtime (Vol. 05)};
  \node[module] (api) at (9.2,0.5) {API Gateway (Vol. 08)};

  % Filesystem / WAL
  \node[draw, rounded corners, dashed, fill=gray!5, text width=10cm, align=center,
        font=\small, inner sep=4pt] (fs) at (4.6,-1.4)
    {Filesystem (WAL): \texttt{.llmwiki/} \{units, graph, index, wal\}};

  % layer background rectangles
  \begin{scope}[on background layer]
    \node[draw, rounded corners, thick, fill=green!10, fit=(parse)(chunk)(embed), inner sep=7pt] {};
    \node[draw, rounded corners, thick, fill=blue!10, fit=(graph)(index)(evidence), inner sep=7pt] {};
    \node[draw, rounded corners, thick, fill=orange!10, fit=(planner)(runtime)(api), inner sep=7pt] {};
  \end{scope}

  % Connections
  \draw[arrow] (parse.east) -- (graph.west);
  \draw[arrow] (chunk.east) -- (index.west);
  \draw[arrow] (embed.east) -- (index.west);
  \draw[arrow] (graph.east) -- (planner.west);
  \draw[arrow] (index.east) -- (planner.west);
  \draw[arrow] (evidence.east) -- (runtime.west);
  \draw[arrow] (planner.east) -- (runtime.west);
  \draw[arrow] (runtime.east) -- (api.west);

  % Filesystem connections
  \draw[arrow] (parse.south) -- (fs.north);
  \draw[arrow] (index.south) -- (fs.north);
\end{tikzpicture}
\end{adjustbox}
\caption{Three-Layer Architecture of LLMWiki. Filesystem acts as write-ahead log (WAL) for all layers.}
\label{fig:architecture:three-layer}
\end{figure}

\subsection{Layer 1: Ingestion Pipeline}
\label{sec:arch:ingestion}
Responsible for parsing heterogeneous artifacts into semantic units - content-addressed, typed, and annotated fragments. Modules:
\begin{itemize}
  \item Semantic Parser (Vol. 02): Language-aware parsing (code, markdown, PDF, OpenAPI, SQL DDL, etc.)
  \item Chunker: Hierarchical, structure-aware segmentation with overlap control
  \item Embedder: Local embedding models (BGE, E5, nomic-embed) with batch inference
\end{itemize}

\subsection{Layer 2: Core Knowledge Space}
\label{sec:arch:core}
The persistent, queryable knowledge representation. Three tightly coupled modules:
\begin{itemize}
  \item Knowledge Graph (Vol. 03): Property graph with typed nodes/edges, entity resolution, and incremental updates
  \item Semantic Index (Vol. 03): Hybrid index (HNSW vector + BM25 keyword + graph adjacency)
  \item Evidence Gate (Vol. 06): Cross-encoder verification + citation extraction + confidence calibration
\end{itemize}

\subsection{Layer 3: Query and Execution}
\label{sec:arch:query}
Composable query execution with planning and observability:
\begin{itemize}
  \item Query Planner (Vol. 07): DAG-based plan compilation, cost-based optimization, subquery decomposition
  \item Agent Runtime (Vol. 05): Tool execution, state management, streaming responses
  \item API Gateway (Vol. 08): REST/gRPC/WebSocket contracts, authentication, rate limiting
\end{itemize}

\section{Filesystem as Write-Ahead Log}
\label{sec:architecture:filesystem-wal}

\begin{designdecision}
\label{dec:fs-as-wal}
\textbf{Decision:} Use the local filesystem (directory structure .llmwiki/) as the write-ahead log instead of a dedicated WAL service (e.g., Kafka, Redis Streams, WAL-G).
\end{designdecision}

\begin{designrationale}
\begin{enumerate}
  \item Zero operational dependency: No separate service to deploy, monitor, or backup.
  \item Deterministic replay: rsync or git of .llmwiki/ fully reproduces state.
  \item Content-addressed deduplication: Units stored as units/<hash>.json are naturally deduplicated.
  \item Atomicity via rename: Write to temp file, rename(2) for atomic commit - POSIX guarantee.
  \item Backup equals copy: tar, rsync, borg, restic work natively.
\end{enumerate}
\end{designrationale}

\begin{designtradeoff}
\begin{itemize}
  \item Concurrency: Single-writer per repository; multi-writer requires external locking (file lock or SQLite WAL mode).
  \item Scale: Directory entry count grows with unit count. Mitigation: sharded subdirectories (units/ab/cd/ef...).
  \item Durability: fsync on commit; configurable \texttt{sync\_on\_write} flag.
\end{itemize}
\end{designtradeoff}

\subsection{Directory Layout}
\label{sec:arch:fs-layout}

\begin{verbatim}
.llmwiki/
|-- units/                    # Content-addressed semantic units
|   |-- ab/cd/ef1234.json     # SHA-256 prefix sharding
|   |-- ...
|-- graph/                    # Knowledge graph persistence
|   |-- nodes/                # Node partitions (by type)
|   |-- edges/                # Edge partitions (by relation type)
|   |-- indexes/              # Graph-local indexes (adjacency, labels)
|-- index/                    # Semantic search indexes
|   |-- hnsw/                 # HNSW vector index (disk-backed)
|   |-- bm25/                 # BM25 inverted index (tantivy)
|   |-- adjacency/            # Graph adjacency for traversal
|-- wal/                      # Write-ahead log (append-only JSONL)
|   |-- 2024-01-15-001.jsonl  # Daily rotation
|   |-- ...
|-- planner/                  # Query plan cache
|   |-- plans/
|-- evidence/                 # Evidence gate caches
|   |-- cross-encoder/
|-- config.toml               # Repository configuration
\end{verbatim}

\section{Data Flow Architecture}
\label{sec:architecture:dataflow}

\begin{figure}[htbp]
\centering
\begin{adjustbox}{max width=\textwidth}
\begin{tikzpicture}[
  stage/.style={draw, rounded corners, fill=blue!5, text width=2.5cm, minimum height=1.2cm, align=center, font=\small, inner sep=3pt},
  rowlabel/.style={font=\bfseries\small, anchor=west},
  flow/.style={->, thick, >=stealth}
]
  % --- Ingestion pipeline: two rows of three ---
  \node[stage] (artifact) {Artifact Input};
  \node[stage, right=0.55cm of artifact] (parse) {Semantic Parse};
  \node[stage, right=0.55cm of parse] (chunk) {Structure-Aware Chunk};

  \node[stage, below=0.9cm of artifact] (extract) {Entity / Relation Extraction};
  \node[stage, right=0.55cm of extract] (embed) {Embed + Index};
  \node[stage, right=0.55cm of embed] (persist) {Persist to FS + Graph};

  \draw[flow] (artifact) -- (parse);
  \draw[flow] (parse) -- (chunk);
  \draw[flow] (chunk.south) -- ++(0,-0.35) -| (extract.north);
  \draw[flow] (extract) -- (embed);
  \draw[flow] (embed) -- (persist);

  % --- Query pipeline: two rows of three ---
  \node[stage, below=1.5cm of extract] (query) {Query Input};
  \node[stage, right=0.55cm of query] (plan) {Query Planner};
  \node[stage, right=0.55cm of plan] (retrieve) {Hybrid Retrieve};

  \node[stage, below=0.9cm of query] (evidence) {Evidence Gate};
  \node[stage, right=0.55cm of evidence] (synthesize) {Synthesize + Cite};
  \node[stage, right=0.55cm of synthesize] (answer) {Answer};

  \draw[flow] (query) -- (plan);
  \draw[flow] (plan) -- (retrieve);
  \draw[flow] (retrieve.south) -- ++(0,-0.35) -| (evidence.north);
  \draw[flow] (evidence) -- (synthesize);
  \draw[flow] (synthesize) -- (answer);

  % Persisted state feeds retrieval
  \draw[flow, dashed] (persist.east) -- ++(0.35,0) |- (retrieve.east);
\end{tikzpicture}
\end{adjustbox}
\caption{End-to-end data flow: ingestion pipeline (top) and query pipeline (bottom).}
\label{fig:architecture:dataflow}
\end{figure}

\subsection{Ingestion Pipeline}
\label{sec:arch:ingestion-pipeline}

\begin{algorithm}[H]
\caption{Ingestion Pipeline}
\label{alg:ingestion}
\begin{algorithmic}[1]
\Require Artifact A (file, API response, database row)
\Ensure Persisted semantic units U, graph updates Delta G, index updates Delta I
\State U = Parse(A) \Comment{Language-aware parser}
\State U = Chunk(U) \Comment{Hierarchical, overlap-controlled}
\State Delta G = ExtractEntitiesRelations(U) \Comment{Local LLM or rule-based}
\State Delta I = Embed(U) \Comment{Batch local inference}
\State WriteWAL(Delta G, Delta I, U) \Comment{Atomic append to WAL}
\State ApplyGraph(Delta G) \Comment{Incremental graph update}
\State ApplyIndex(Delta I) \Comment{Incremental index update}
\State \Return U, Delta G, Delta I
\end{algorithmic}
\end{algorithm}

\subsection{Query Pipeline}
\label{sec:arch:query-pipeline}

\begin{algorithm}[H]
\caption{Query Execution Pipeline}
\label{alg:query}
\begin{algorithmic}[1]
\Require Query q, context budget B, evidence threshold tau
\Ensure Answer a with citations C
\State $P$ = Plan($q$) \Comment{Planner outputs DAG of subqueries}
\State $R = \emptyset$
\ForAll{subquery $s$ in $P$}
  \State $R_s$ = HybridRetrieve($s$) \Comment{Vector + BM25 + Graph}
  \State $R = R \cup R_s$
\EndFor
\State E = EvidenceGate(R, tau) \Comment{Cross-encoder verify}
\State a = Synthesize(E, B) \Comment{LLM with citations forced}
\State \Return a, C(E) \Comment{Answer + evidence citations}
\end{algorithmic}
\end{algorithm}

\section{Core Invariants}
\label{sec:architecture:invariants}

\begin{invariant}[Content Addressability]
\label{inv:content-address}
Forall u in U: id(u) = SHA256(canonical(u)). Identical content implies identical identifier across all indices.
\end{invariant}

\begin{invariant}[Evidence Closure]
\label{inv:evidence-closure}
Forall a in Answers: citations(a) subset of U and forall c in citations(a): verify(c) >= tau.
\end{invariant}

\begin{invariant}[Incremental Consistency]
\label{inv:incremental}
After any ingestion I, the triple (G, I, U) satisfies: query(q, G, I, U) = replay(queries(q)).
\end{invariant}

\begin{invariant}[Local-First Determinism]
\label{inv:local-determinism}
Given identical .llmwiki/ directory and identical query q, two executions produce bitwise-identical results (modulo non-deterministic LLM sampling, which is confined to planning/synthesis).
\end{invariant}

\section{The Filesystem as Lineage Recorder}
\label{sec:architecture:lineage}

Invariants~\ref{inv:incremental} and \ref{inv:local-determinism} assert that history can be replayed. They do not say how history is \emph{reconstructed} when the record is incomplete, and every long-lived deployment produces an incomplete record: WAL segments are compacted, old snapshots are pruned, repositories are cloned from a mid-life state. A replay guarantee that holds only over an unbroken log is a guarantee that expires.

Developmental biology has spent a decade on exactly this reconstruction problem. CRISPR lineage recorders accumulate irreversible edits at engineered target sites across cell divisions, so that the division history of an organism can be inferred from the barcodes of its terminal cells alone --- the intermediate states are gone, and the tree is recovered from the endpoint. Recent surveys note that recorders which log a \emph{deterministic temporal ordering} of mutations enable precise reconstruction of deep lineages \citep{lineagetracing2026}, which is precisely the argument for an append-only WAL. More usefully, the field has formalised inference under a degraded record: Bayesian phylodynamic frameworks reconstruct lineage trees with calibrated uncertainty from barcode data with dependent target sites and missing observations \citep{phylodynamic2025}.

Two failure modes from that literature transfer to LLMWiki without modification.

\paragraph{Homoplasy.} An identical edit arising independently in two lineages falsely implies shared ancestry. Content addressability (Invariant~\ref{inv:content-address}) makes this endemic rather than occasional: two units produced by unrelated derivation paths that happen to canonicalise identically receive the same identifier \emph{by design}. This is the deduplication property, and it is desirable for storage and indexing. It is also destructive of provenance, because it collapses distinct derivations into one node. Identity and ancestry are different relations and must not share a key.

\begin{invariant}[Lineage Distinctness]
\label{inv:lineage-distinct}
Content identity does not imply derivation identity. For a unit $u$, $\text{id}(u)$ addresses content; the derivation of $u$ is recorded separately as the set of WAL edges producing it. Provenance is therefore a directed acyclic graph over content identifiers, never a tree, and a query for the origin of $u$ returns a set of derivations rather than a unique parent.
\end{invariant}

\paragraph{Dropout.} A target site that fails to be read yields a missing character, not a wrong one, and treating the two alike corrupts the tree. The analogue is a compacted WAL segment: the derivation edge is unobserved, not absent. Recording compaction boundaries explicitly lets reconstruction distinguish \emph{no derivation occurred} from \emph{the derivation is unobservable}, which is the difference between a false claim and an honest one.

\begin{invariant}[Replay Degradation is Explicit]
\label{inv:replay-degradation}
Where the WAL has been compacted, replay returns a result annotated with the reconstruction's uncertainty rather than an unqualified answer. A deployment that cannot verify a replay reports that it cannot, and never silently substitutes a plausible history.
\end{invariant}

This is the mechanism that makes the auditability claim survive contact with operations. A system that guarantees reproducibility only while nothing has been pruned has guaranteed reproducibility only in the laboratory; the useful statement is that reproducibility degrades in a declared, measurable way.

\section{Comparison with Prior Art}
\label{sec:architecture:prior-art}

\begin{table}[htbp]
\centering
\caption{LLMWiki vs. Prior Art: Architectural Positioning}
\label{tab:arch:prior-art}
\begin{adjustbox}{max width=\textwidth}
\begin{tabular}{lcccc}
\toprule
\textbf{Property} & \textbf{RAG} & \textbf{GraphRAG} & \textbf{HippoRAG} & \textbf{LLMWiki} \\
\midrule
Graph as source of truth & No & Yes & Yes & Yes \\
Vectors as projection & No & No & No & Yes \\
Filesystem as WAL & No & No & No & Yes \\
Evidence gate (first-class) & No & No & No & Yes \\
Local-first (no API required) & No & No & No & Yes \\
Deterministic replay & No & No & No & Yes \\
Content-addressed units & No & No & No & Yes \\
Query planner (DAG) & No & No & No & Yes \\
\bottomrule
\end{tabular}
\end{adjustbox}
\end{table}

\section{Patent-Relevant Novelty}
\label{sec:architecture:patent}

\begin{patentclaim}
\textbf{Claim 1:} A semantic operating layer characterized by: (a) content-addressed semantic units stored in a filesystem-backed write-ahead log, (b) a property knowledge graph as the canonical knowledge representation, (c) a hybrid semantic index combining vector, keyword, and graph adjacency, (d) an evidence gate that verifies retrieval results via cross-encoder scoring before synthesis, and (e) a DAG-based query planner that decomposes queries into verifiable subqueries.
\end{patentclaim}

\begin{patentclaim}
\textbf{Claim 2:} The system of Claim 1 wherein the filesystem serves as the sole durability mechanism, enabling deterministic replay via directory synchronization without auxiliary transaction logs.
\end{patentclaim}

\begin{patentclaim}
\textbf{Claim 3:} The system of Claim 1 wherein the evidence gate operates as a mandatory pipeline stage, rejecting any retrieval result that fails cross-encoder verification against a calibrated threshold, thereby enforcing evidence closure.
\end{patentclaim}

\begin{patentdiscussion}
Prior art distinctions:
\begin{itemize}
  \item RAG (Lewis et al., 2020): Vector index is primary; no graph; no evidence gate; API-dependent.
  \item GraphRAG (Edge et al., 2024): Graph constructed via LLM; no content-addressed units; no filesystem WAL; no mandatory evidence gate.
  \item HippoRAG (Zhang et al., 2024): Personalized PageRank over graph; no hybrid index; no deterministic replay.
  \item LightRAG (Zeng et al., 2024): Dual-level retrieval; no evidence gate; no DAG planner.
  \item RAPTOR (Sarthi et al., 2024): Recursive summarization tree; no graph; no filesystem WAL.
\end{itemize}
\end{patentdiscussion}

\section{Module Interface Contracts}
\label{sec:architecture:contracts}

\begin{interface}[Parser]
\label{iface:parser}
\begin{verbatim}
trait Parser {
    fn parse(&self, artifact: Artifact) -> Result<Vec<SemanticUnit>>;
    fn supported_mimes(&self) -> Vec<MimeType>;
    fn schema(&self) -> ParserSchema;
}
\end{verbatim}
\end{interface}

\begin{interface}[Chunker]
\label{iface:chunker}
\begin{verbatim}
trait Chunker {
    fn chunk(&self, unit: SemanticUnit, config: ChunkConfig) -> Vec<Chunk>;
    fn config_schema(&self) -> ChunkConfigSchema;
}
\end{verbatim}
\end{interface}

\begin{interface}[Embedder]
\label{iface:embedder}
\begin{verbatim}
trait Embedder {
    fn embed_batch(&self, texts: &[String]) -> Result<Vec<Vec<f32>>>;
    fn dimension(&self) -> usize;
    fn model_id(&self) -> ModelId;
}
\end{verbatim}
\end{interface}

\begin{interface}[GraphStore]
\label{iface:graphstore}
\begin{verbatim}
trait GraphStore {
    fn upsert_nodes(&self, nodes: Vec<Node>) -> Result<()>;
    fn upsert_edges(&self, edges: Vec<Edge>) -> Result<()>;
    fn traverse(&self, start: NodeId, pattern: TraversalPattern) -> Result<Vec<Path>>;
    fn neighborhoods(&self, ids: &[NodeId], hops: u32) -> Result<Subgraph>;
}
\end{verbatim}
\end{interface}

\begin{interface}[SemanticIndex]
\label{iface:semantic-index}
\begin{verbatim}
trait SemanticIndex {
    fn upsert(&self, units: Vec<IndexedUnit>) -> Result<()>;
    fn search_vector(&self, query: Vec<f32>, k: usize) -> Result<Vec<ScoredUnit>>;
    fn search_bm25(&self, query: &str, k: usize) -> Result<Vec<ScoredUnit>>;
    fn search_hybrid(&self, query: HybridQuery, k: usize) -> Result<Vec<ScoredUnit>>;
}
\end{verbatim}
\end{interface}

\begin{interface}[EvidenceGate]
\label{iface:evidence-gate}
\begin{verbatim}
trait EvidenceGate {
    fn verify(&self, query: &str, candidates: Vec<ScoredUnit>, threshold: f32) 
        -> Result<Vec<VerifiedUnit>>;
    fn calibrate(&self, dataset: CalibrationSet) -> Result<CalibrationCurve>;
}
\end{verbatim}
\end{interface}

\begin{interface}[QueryPlanner]
\label{iface:query-planner}
\begin{verbatim}
trait QueryPlanner {
    fn plan(&self, query: Query, context: PlanContext) -> Result<QueryPlan>;
    fn optimize(&self, plan: QueryPlan, stats: IndexStats) -> Result<QueryPlan>;
    fn explain(&self, plan: QueryPlan) -> PlanExplanation;
}
\end{verbatim}
\end{interface}

\section{Runtime Configuration}
\label{sec:architecture:config}

\begin{verbatim}
# .llmwiki/config.toml
[repository]
id = "llmwiki-main"
created_at = "2024-01-15T10:30:00Z"
version = "0.1.0"

[ingestion]
parsers = ["python", "markdown", "pdf", "openapi", "sql", "yaml", "json"]
chunker = { type = "ast-aware", max_tokens = 512, overlap = 50 }
embedder = { model = "bge-m3", batch_size = 32, device = "cpu" }

[graph]
store = "rocksdb"
entity_resolution = { enabled = true, threshold = 0.85 }
index_types = ["btree", "adjacency", "label"]

[index]
hnsw = { m = 16, ef_construction = 200, ef_search = 64 }
bm25 = { k1 = 1.2, b = 0.75 }
hybrid = { vector_weight = 0.6, keyword_weight = 0.3, graph_weight = 0.1 }

[evidence_gate]
model = "cross-encoder/ms-marco-MiniLM-L-6-v2"
threshold = 0.72
calibration = "isotonic"
batch_size = 16

[planner]
max_depth = 5
max_branching = 4
cost_model = "learned"
cache_ttl = "1h"

[api]
bind = "127.0.0.1:8080"
auth = "none"
rate_limit = 100
\end{verbatim}

\section{Extension Points}
\label{sec:architecture:extensions}

\begin{table}[htbp]
\centering
\caption{LLMWiki Extension Points}
\label{tab:arch:extensions}
\begin{adjustbox}{max width=\textwidth}
\begin{tabular}{lll}
\toprule
\textbf{Extension Point} & \textbf{Interface} & \textbf{Volume} \\
\midrule
Custom Parser & Parser & Vol. 02 \\
Custom Chunker & Chunker & Vol. 02 \\
Custom Embedder & Embedder & Vol. 02 \\
Graph Store Backend & GraphStore & Vol. 03 \\
Index Backend & SemanticIndex & Vol. 03 \\
Evidence Verifier & EvidenceGate & Vol. 06 \\
Query Planner Strategy & QueryPlanner & Vol. 07 \\
Agent Tool & Tool & Vol. 05 \\
Auth Provider & AuthProvider & Vol. 08 \\
\bottomrule
\end{tabular}
\end{adjustbox}
\end{table}

\section{Summary}
\label{sec:architecture:summary}

LLMWiki's architecture is defined by five differentiating choices:
\begin{enumerate}
  \item \textbf{Graph-first, vectors-second}: The property graph is the canonical knowledge representation; vectors are derived projections for similarity search.
  \item \textbf{Filesystem as WAL}: Zero operational dependencies; deterministic replay via directory sync.
  \item \textbf{Content-addressed units}: SHA-256 identities enable deduplication, caching, and verifiable lineage.
  \item \textbf{Evidence gate as mandatory stage}: No retrieval result reaches synthesis without cross-encoder verification.
  \item \textbf{DAG query planner}: Queries decompose into verifiable subqueries with cost-based optimization.
\end{enumerate}
These choices are not incremental improvements over RAG - they represent a different architectural category: a semantic operating layer.

The subsequent volumes detail each module's algorithms, data structures, APIs, and implementation considerations.

%===============================================================================
% End of Volume 01
%===============================================================================
